%% 自動生成（tools/mkchecklist.py）。各章の「ふりかえり」の項目を集めたもの
\begingroup\footnotesize
\def\bkcb{\tikz[baseline=-0.5ex]\draw[\bkTF, line width=0.7pt, rounded corners=0.5pt] (0,-0.3em) rectangle (0.8em,0.5em);}
\begin{longtblr}{colspec={Q[c,m,wd=9mm] X[l,m] Q[c,m,wd=10mm] Q[c,m,wd=9mm] Q[c,m,wd=9mm]}, width=\linewidth, rowhead=1, colsep=3pt,
  row{1}={bg=bkVol!14, font=\sffamily\bfseries\scriptsize}, hlines={0.4pt,bkLine}, rowsep=0.7pt}
章 & できるようになったこと & ページ & 1回め & 2回め \\
{\sffamily\bfseries\color{\bkTT}4-1} & 平行線の同位角・錯角から、角を求められる。 & \pageref{bk:4-1:heikou} & \bkcb & \bkcb \\
 & 内角の和が180°のわけと、多角形の角の和がわかる。 & \pageref{bk:4-1:naikaku} & \bkcb & \bkcb \\
 & 定義・公理・定理・法則のちがいを説明でき、合同条件・逆・反例を使える。 & \pageref{bk:4-1:teigi} & \bkcb & \bkcb \\
{\sffamily\bfseries\color{\bkTT}4-2} & 弦・弧・おうぎ形がわかり、弧の長さとおうぎ形の面積を、中心角に比例させて求められる。 & \pageref{bk:4-2:hirei} & \bkcb & \bkcb \\
 & 角をラジアンで表し、度と換算して、$s=r\theta$ と $S=\tfrac{1}{2}r^2\theta$ を使える。 & \pageref{bk:4-2:rad} & \bkcb & \bkcb \\
 & 接線が半径に垂直なわけを説明でき、DEG・RAD と極座標 $(r,\ \theta)$ を使える。 & \pageref{bk:4-2:sessen} & \bkcb & \bkcb \\
{\sffamily\bfseries\color{\bkTT}4-3} & 拡大図・縮図と縮尺を使って、本当の長さや図の上の長さを求められる。 & \pageref{bk:4-3:kakudai} & \bkcb & \bkcb \\
 & 三角形の相似条件を見つけ、相似比や平行線と線分の比から、辺の長さを求められる。 & \pageref{bk:4-3:jouken} & \bkcb & \bkcb \\
 & 面積の比は $k^2$、体積の比は $k^3$ を使い、すいの断面や直角三角形の辺の比を説明できる。 & \pageref{bk:4-3:menseki} & \bkcb & \bkcb \\
{\sffamily\bfseries\color{\bkTT}4-4} & 方眼の面積から $a^2+b^2=c^2$ を見つけ、2通りに説明できる。 & \pageref{bk:4-4:shoumei} & \bkcb & \bkcb \\
 & 辺の長さ・2点の距離・直方体の対角線を求められる。 & \pageref{bk:4-4:kyori} & \bkcb & \bkcb \\
 & 直角三角形かを見分け、円の式のわけがわかる。 & \pageref{bk:4-4:en} & \bkcb & \bkcb \\
{\sffamily\bfseries\color{\bkTT}4-5} & 平行四辺形・三角形・台形・ひし形の面積を、長方形に直して求められる。 & \pageref{bk:4-5:sankaku} & \bkcb & \bkcb \\
 & 円の面積と、柱・すい・球の体積を、わけといっしょに使える。 & \pageref{bk:4-5:en} & \bkcb & \bkcb \\
 & 展開図から表面積を求め、体積を L に換算できる。 & \pageref{bk:4-5:tenkai} & \bkcb & \bkcb \\
{\sffamily\bfseries\color{\bkTT}4-6} & 直角三角形で、どの角から見るかに気をつけて、$\sin$・$\cos$・$\tan$ を求められる。 & \pageref{bk:4-6:teigi} & \bkcb & \bkcb \\
 & 三角比と電卓で、辺の長さや角度を求められる。坂の傾きを角度に直せる。 & \pageref{bk:4-6:atai} & \bkcb & \bkcb \\
 & 斜面の重力を、$mg\sin\theta$ と $mg\cos\theta$ に分けられる。 & \pageref{bk:4-6:bunkai} & \bkcb & \bkcb \\
{\sffamily\bfseries\color{\bkTT}4-7} & 鈍角の三角比を使い、三角形の面積を $\tfrac{1}{2}ab\sin C$ で求められる。 & \pageref{bk:4-7:menseki} & \bkcb & \bkcb \\
 & 正弦定理と余弦定理を使い分けて、三角形の辺や角を求められる。 & \pageref{bk:4-7:seigen} & \bkcb & \bkcb \\
 & ロボットのうでの関節の角を、余弦定理で求められる。 & \pageref{bk:4-7:robot} & \bkcb & \bkcb \\
{\sffamily\bfseries\color{\bkTT}4-8} & 単位円の上の点 $(\cos\theta,\ \sin\theta)$ を求められる。 & \pageref{bk:4-8:tanien} & \bkcb & \bkcb \\
 & $\sin^2\theta+\cos^2\theta=1$ や90°回す式を説明できる。 & \pageref{bk:4-8:sanheihou} & \bkcb & \bkcb \\
 & $\omega$ と $T$ で、回る点の位置・速さを求められる。 & \pageref{bk:4-8:omega} & \bkcb & \bkcb \\
{\sffamily\bfseries\color{\bkTT}4-9} & 加法定理を単位円で導き、$\sin(a\pm b)$・$\cos(a\pm b)$ を計算できる。 & \pageref{bk:4-9:kahou} & \bkcb & \bkcb \\
 & 2倍角の式と $\sin^2\theta=\tfrac{1-\cos2\theta}{2}$ を作り、使える。 & \pageref{bk:4-9:nibai} & \bkcb & \bkcb \\
 & 合成と「積と和の式」で、うなりや電力を説明できる。 & \pageref{bk:4-9:gousei} & \bkcb & \bkcb \\
{\sffamily\bfseries\color{\bkTT}4-10} & 向きと大きさをもつ量を矢印と記号で表せる。 & \pageref{bk:4-10:kigou} & \bkcb & \bkcb \\
 & 矢印をつないで、合力の大きさを求められる。 & \pageref{bk:4-10:wa} & \bkcb & \bkcb \\
 & 差をかいて、速度の変化や相対速度を求められる。 & \pageref{bk:4-10:sa} & \bkcb & \bkcb \\
{\sffamily\bfseries\color{\bkTT}4-11} & 成分で、和・差・実数倍・大きさを計算できる。 & \pageref{bk:4-11:keisan} & \bkcb & \bkcb \\
 & 力や速度を、成分に分けて考えられる。 & \pageref{bk:4-11:bunkai} & \bkcb & \bkcb \\
 & 内分点・重心・進む点の位置を求められる。 & \pageref{bk:4-11:ichi} & \bkcb & \bkcb \\
{\sffamily\bfseries\color{\bkTT}4-12} & 仕事を求め、$W=Fs\cos\theta$ を使える。 & \pageref{bk:4-12:naname} & \bkcb & \bkcb \\
 & 内積を、大きさと角からも、成分からも計算できる。 & \pageref{bk:4-12:naiseki} & \bkcb & \bkcb \\
 & 内積で、なす角・直交・相関係数を求められる。 & \pageref{bk:4-12:chokkou} & \bkcb & \bkcb \\
{\sffamily\bfseries\color{\bkTT}4-13} & 空間の点を $(x,\ y,\ z)$ で表し、2点の距離を求められる。 & \pageref{bk:4-13:zahyo} & \bkcb & \bkcb \\
 & 3次元のベクトルを成分で計算し、内積からなす角を求められる。 & \pageref{bk:4-13:naiseki} & \bkcb & \bkcb \\
 & 基本ベクトル・単位ベクトル・法線ベクトルで、向きや平面を表せる。 & \pageref{bk:4-13:kihon} & \bkcb & \bkcb \\
{\sffamily\bfseries\color{\bkTT}4-14} & 力のモーメントを「力×うでの長さ」で求め、$M=xF_y-yF_x$ で計算できる。 & \pageref{bk:4-14:moment} & \bkcb & \bkcb \\
 & 外積の向きと大きさがわかり、成分で計算できる。 & \pageref{bk:4-14:gaiseki} & \bkcb & \bkcb \\
 & トルクや、磁場から受ける力の向きを、外積で求められる。 & \pageref{bk:4-14:torque} & \bkcb & \bkcb \\
{\sffamily\bfseries\color{\bkTT}4-15} & 極座標と $(x,\ y)$ を、区画に注意して行き来できる。 & \pageref{bk:4-15:kyoku} & \bkcb & \bkcb \\
 & $\vec e_r$、$\vec e_\theta$ を使って、円運動の速度を $r\omega\,\vec e_\theta$ と表せる。 & \pageref{bk:4-15:er} & \bkcb & \bkcb \\
 & 円柱座標・球座標で空間の点を表し、$(x,\ y,\ z)$ に直せる。 & \pageref{bk:4-15:enchu} & \bkcb & \bkcb \\
\end{longtblr}
\endgroup
